\chapter{Marco teórico y revisión de la literatura}
\label{cap:marco-teorico}

% Responde a ¿cómo? y ¿con qué herramientas?: define los conceptos clave del trabajo.

Este capítulo define los conceptos y las herramientas sobre los que se apoya el
resto de la memoria, siguiendo una lógica deductiva que va de las nociones
generales a los instrumentos concretos. Primero se precisan las ideas de agenda
temática, \emph{salience} y coherencia partidaria, y la distinción entre lo que
un partido declara y lo que revela al votar, que es la pregunta que articula el
trabajo. Luego se presenta la taxonomía del \emph{Manifesto Project} (MARPOR)
como espacio de categorías políticas comparable entre fuentes, junto con el índice
posicional RILE y las encuestas de expertos que sirven de referencia externa.
Por último se revisan las técnicas de procesamiento de lenguaje natural que
permiten clasificar texto político a gran escala. Estos conceptos preparan la
metodología del capítulo~\ref{cap:metodologia} y la lectura de los hallazgos del
capítulo~\ref{cap:resultados}, pero aquí no se presentan resultados del caso
estudiado.

\section{Discurso político, agenda y coherencia partidaria}
\label{sec:discurso-coherencia}

El estudio del comportamiento partidario distingue de manera habitual entre lo
que un partido \emph{dice} y lo que un partido \emph{hace}. La primera dimensión
remite a la \textbf{agenda declarada}, es decir, al conjunto de temas que una
organización política decide poner en primer plano en su comunicación pública, ya
sea en su programa electoral, en los medios o en el debate parlamentario. La
segunda remite a la \textbf{agenda revelada}, esto es, a las posiciones que el
partido expresa cuando sus representantes votan sobre textos concretos. Ambas son
observables, pero no son equivalentes, y su contraste es el problema central de
esta memoria.

Un concepto útil para caracterizar la agenda declarada es el de \emph{salience},
la prominencia relativa que un actor otorga a cada tema. La literatura sobre
competencia partidaria sostiene que los partidos compiten menos tomando
posiciones opuestas sobre los mismos asuntos y más enfatizando selectivamente
aquellos temas en los que se sienten fuertes, una idea vinculada a la noción de
\emph{issue ownership}, la asociación duradera entre un partido y ciertos temas.
Desde esta perspectiva, describir una agenda es, ante todo, medir de qué se habla
y cuánto, antes que ubicar al actor en una escala de posiciones.

Esta óptica temática tiene una consecuencia metodológica importante. Reducir a un
partido a un único punto en un eje izquierda--derecha resume su posición, pero
pierde la estructura de su agenda y resulta frágil cuando el objeto de interés es
la coherencia entre discurso y voto. La ciencia política comparada ha mostrado,
además, que la competencia europea no se ordena solo por el eje económico, sino
también por una segunda dimensión que opone valores libertarios, alternativos y
cosmopolitas a posiciones tradicionales, autoritarias y nacionalistas, conocida
como dimensión GAL--TAN \cite{hooghe2002galtan}. Un enfoque basado en el énfasis
por temas permite captar esa multidimensionalidad, que un índice unidimensional
tiende a aplanar.

La agenda declarada, sin embargo, no se expresa igual en todas partes. Un mismo
partido se dirige a audiencias distintas en su programa, en las redes sociales y
en el hemiciclo, y adapta a cada canal tanto el registro como la selección de
temas, de modo que el canal opera como un condicionante estructural de lo que se
observa \cite{ivanusch2024channels}. Comparar la agenda de un partido consigo
misma a través de canales, y no solo entre partidos, se vuelve entonces una
pregunta con contenido propio.

Del lado de la agenda revelada, el estudio del voto legislativo aporta un segundo
cuerpo de conceptos. En los sistemas parlamentarios, el voto rara vez es la suma
de decisiones individuales independientes: los grupos exhiben una fuerte
\textbf{disciplina de bloque}, sostenida por incentivos organizativos y por la
lógica de gobierno frente a oposición \cite{sieberer2006partyunity}. La unidad de
voto observada es, en consecuencia, el resultado de principales en competencia,
como el partido, el electorado y el gobierno, más que un reflejo directo de las
preferencias temáticas de cada representante \cite{carey2007competing}. Esta
distinción anticipa una idea que la memoria retoma más adelante: el voto puede
revelar sobre todo una posición de bloque y, solo en ciertos márgenes, una señal
temática propia. El contraste sistemático entre la agenda declarada y la revelada,
sobre un mismo conjunto de actores y bajo una taxonomía común, es precisamente el
vacío que este trabajo busca ocupar.

\section{El Manifesto Project y la taxonomía MARPOR}
\label{sec:marpor}

Para comparar el contenido político de fuentes tan distintas como un programa
electoral, un mensaje en redes o un texto legislativo se necesita un espacio de
categorías común. Esta memoria adopta el esquema del \emph{Manifesto Project}
(MARPOR) \cite{marpor_mp_v5}, un estándar consolidado de la ciencia política
comparada con el que se han codificado manualmente programas electorales de
numerosos países durante décadas. Su unidad de análisis es la \emph{cuasi-frase},
un fragmento que expresa una única idea política, a la que se asigna una
categoría temática.

El esquema organiza el contenido en \textbf{56 categorías} agrupadas en
\textbf{siete dominios} macro. El dominio se obtiene del primer dígito del código
de categoría, lo que habilita una lectura jerárquica: fina a nivel de categoría y
más robusta a nivel de dominio. El Cuadro~\ref{tab:marpor-dominios} resume los
siete dominios con una categoría de ejemplo cada uno; el listado completo de las
56 categorías se recoge en el anexo~\ref{anexo:marpor}.

\begin{table}[H]
\centering
\caption{Los siete dominios de la taxonomía MARPOR, con una categoría ilustrativa
por dominio. El detalle de las 56 categorías se encuentra en el
anexo~\ref{anexo:marpor}.}
\label{tab:marpor-dominios}
\small
\begin{tabularx}{\textwidth}{@{}l l >{\raggedright\arraybackslash}X@{}}
\toprule
\textbf{Dominio} & \textbf{Nombre (MARPOR)} & \textbf{Categoría de ejemplo} \\
\midrule
1 & External Relations & Internacionalismo \\
2 & Freedom and Democracy & Libertad y derechos humanos \\
3 & Political System & Autoridad política \\
4 & Economy & Economía de libre mercado \\
5 & Welfare and Quality of Life & Expansión del Estado de bienestar \\
6 & Fabric of Society & Ley y orden \\
7 & Social Groups & Grupos laborales \\
\bottomrule
\end{tabularx}
\end{table}

Tres propiedades hacen de MARPOR un espacio idóneo para este trabajo. Es
\emph{independiente de la fuente}, porque las mismas categorías se aplican por
igual a distintos tipos de texto; es \emph{citable y externamente validable}, al
ser un estándar internacional cuyas estimaciones pueden contrastarse contra
codificaciones humanas y encuestas de expertos; y admite una \emph{lectura tanto
temática como posicional}, pues las distribuciones de categorías sostienen el
análisis de énfasis y, al mismo tiempo, permiten derivar un índice de posición
cuando se lo necesita.

Ese índice es el \textbf{RILE} (\emph{right--left}), propuesto por Laver y Budge
\cite{laver1992party}, que resume la posición de un partido como la diferencia
entre el porcentaje de cuasi-frases asignadas a un conjunto canónico de
categorías de derecha y el de un conjunto de categorías de izquierda. El RILE es
un resumen conveniente, pero tiene límites conceptuales conocidos: es
\emph{unidimensional}, de modo que colapsa en un solo eje una agenda que suele ser
multidimensional, y es \emph{ciego a la dirección} del enunciado, ya que cuenta
del mismo modo el apoyo y el ataque a una misma idea. Por estas razones, en esta
memoria el RILE no es el instrumento del análisis sustantivo, sino que se reserva
para \emph{validar} el pipeline; el análisis principal se construye sobre el
\textbf{énfasis por dominio}, es decir, sobre las distribuciones temáticas
completas y no sobre un índice posicional colapsado.

La validación posicional se apoya en una referencia externa e independiente del
texto: el \emph{Chapel Hill Expert Survey} (CHES), una encuesta que sitúa a los
partidos europeos en distintas dimensiones a partir del juicio de expertos. Se
emplea la edición de 2019 \cite{jolly2022ches, ches2019data}; existe una edición
más reciente de 2024 \cite{rovny2025ches}, que no se utiliza aquí. CHES cumple en
este trabajo un papel estrictamente de \emph{benchmark} externo con el cual
contrastar las posiciones estimadas, y no el de una taxonomía que oriente el
análisis temático.

\section{Procesamiento de lenguaje natural para texto político}
\label{sec:nlp-politico}

Codificar manualmente cinco corpus con el esquema MARPOR sería inviable por su
volumen, de modo que la asignación de categorías se automatiza con técnicas de
procesamiento de lenguaje natural. Se combinan dos enfoques complementarios: una
clasificación supervisada, que asigna las categorías teóricas del esquema, y un
modelado de tópicos no supervisado, que descubre agrupaciones empíricas dentro de
cada corpus. Ambos se describen a continuación en el plano conceptual; su
configuración concreta se detalla en el capítulo~\ref{cap:metodologia}.

\subsection{Clasificación supervisada: ManifestoBERTa}
\label{sec:manifestoberta}

La clasificación supervisada parte de un modelo de lenguaje preentrenado que se
afina para una tarea específica. El núcleo del análisis emplea ManifestoBERTa
\cite{manifestoberta2024}, un modelo publicado por el propio Manifesto Project que
clasifica cuasi-frases en las 56 categorías del esquema MARPOR. Está construido
sobre XLM-RoBERTa \cite{conneau2020xlmr}, un modelo multilingüe preentrenado sobre
texto de numerosos idiomas, lo que resulta apropiado para un corpus en francés y
permite aprovechar representaciones lingüísticas transferibles entre lenguas.

Un supuesto conceptual debe hacerse explícito. ManifestoBERTa fue entrenado sobre
programas electorales, un género con un registro particular, de manera que
aplicarlo a otras superficies, como mensajes en redes, intervenciones
parlamentarias o textos legislativos, constituye una \textbf{transferencia de
dominio}. Esa transferencia es razonable, porque la taxonomía es temática y no
depende del formato, pero no es gratuita, y por eso el desempeño del clasificador
fuera de su dominio de origen debe evaluarse antes de interpretar cualquier
resultado. La forma de esa evaluación se desarrolla en la
sección~\ref{sec:metodo-manifestoberta}.

\subsection{Modelado de tópicos no supervisado: BERTopic}
\label{sec:bertopic}

Como complemento exploratorio se utiliza BERTopic \cite{grootendorst2022bertopic},
una técnica de modelado de tópicos que, en lugar de imponer categorías teóricas,
descubre agrupaciones de documentos a partir de su similitud semántica. Su lógica
encadena varias etapas: representa cada documento como un vector denso mediante
\emph{embeddings} de oraciones \cite{reimers2019sentencebert}, reduce la
dimensionalidad de esos vectores con UMAP \cite{mcinnes2018umap}, agrupa los
documentos con un algoritmo de \emph{clustering} basado en densidad, HDBSCAN
\cite{mcinnes2017hdbscan}, y describe cada grupo con sus términos más
característicos.

El papel de BERTopic en esta memoria es \textbf{exploratorio y no comparable entre
arenas}. Como los tópicos emergen de cada corpus por separado, no constituyen un
espacio común entre fuentes ni permiten contrastar partidos de manera sistemática,
a diferencia de la clasificación supervisada sobre MARPOR. Por ello se emplea para
inspeccionar qué temas afloran en el texto y para detectar ruido o lenguaje
procedimental, pero no para sostener conclusiones cuantitativas. Su configuración
se detalla en la sección~\ref{sec:metodo-bertopic}. El
Cuadro~\ref{tab:nlp-comparacion} resume la complementariedad conceptual de ambos
enfoques.

\begin{table}[H]
\centering
\caption{Comparación conceptual de los dos enfoques de clasificación de texto
empleados en el trabajo.}
\label{tab:nlp-comparacion}
\small
\begin{tabularx}{\textwidth}{@{}l >{\raggedright\arraybackslash}X >{\raggedright\arraybackslash}X@{}}
\toprule
 & \textbf{ManifestoBERTa (supervisado)} & \textbf{BERTopic (no supervisado)} \\
\midrule
Categorías & Fijas y teóricas (esquema MARPOR) & Emergentes del propio corpus \\
Comparabilidad entre arenas & Alta (espacio común) & Baja (tópicos por corpus) \\
Rol en el trabajo & Núcleo del análisis & Complemento exploratorio \\
\bottomrule
\end{tabularx}
\end{table}

\subsection{Detección de postura y marcos de valores}
\label{sec:stance-valores}

Existen enfoques que apuntan a dimensiones del texto político distintas del tema.
La \emph{detección de postura} busca determinar si un texto está a favor o en
contra de un objetivo, y se ha desarrollado, por ejemplo, sobre plataformas de
democracia participativa y consultas públicas europeas
\cite{barriere2022cofe, barriere2023multilingual}. Los \emph{marcos de valores}
ofrecen otra lente complementaria, ya sea a través de la teoría de valores básicos
de Schwartz \cite{schwartz2012overview} o de la teoría de los fundamentos morales
\cite{graham2013moral}, que caracterizan el discurso según los valores o
principios que moviliza. Estos enfoques son afines al problema de esta memoria,
pero quedan \textbf{fuera de su alcance}: el análisis se concentra en el énfasis
temático bajo la taxonomía MARPOR, y la incorporación de la postura o de los
marcos de valores se deja señalada como línea de trabajo futuro.

\section{Trabajos relacionados}
\label{sec:trabajos-relacionados}

La pregunta por la coherencia entre lo que los partidos declaran y lo que hacen
tiene una larga tradición en la ciencia política, que ha estudiado por separado
tanto el contenido de los programas electorales como el comportamiento de voto en
las legislaturas. La codificación de manifiestos y los índices posicionales
derivados de ella han permitido caracterizar la oferta programática de los
partidos, mientras que el estudio del voto legislativo ha documentado la fuerza de
la disciplina de bloque y la lógica de gobierno frente a oposición. Con la
disponibilidad creciente de datos parlamentarios abiertos y de comunicación
digital, ha surgido además una línea que compara la agenda de los partidos a
través de distintos canales de comunicación \cite{ivanusch2024channels}.

El aporte de este trabajo se sitúa en la intersección de esas líneas. En lugar de
analizar una sola fuente, integra cinco arenas discursivas heterogéneas bajo una
cohorte común de diputados y las proyecta sobre una misma taxonomía temática, lo
que permite contrastar de forma directa la agenda declarada con la revelada por el
voto sobre un mismo conjunto de actores. Ese diseño, más que un resultado
sustantivo puntual, es la contribución que la memoria desarrolla en detalle a
partir del capítulo~\ref{cap:metodologia}, donde se describen las fuentes, la
construcción de los corpus y los métodos de clasificación y de análisis.
